Let $b = 2r$; thus $1 \leq b \leq a$.
Square both sides of the equation:
\begin{equation*}
\paren{1 - \frac{1}{a} \cos^2\paren{x}}^{b} = \sin^2\paren{x}
\end{equation*}
Make the substitution $y = 1 - \mfrac{1}{a} \cos^2\paren{x}$.
The range of possibe values of $y$ is $\cinterval{1 - \mfrac{1}{a}}{1}$.
\newline
Using that $\sin^2\paren{x} = 1 - \cos^2\paren{x}$ and $\cos^2\paren{x} = a - a y$,
the substitution yields the equation
\begin{equation*}
y^b - a y + a - 1 = 0
\end{equation*}
Clearly, $1$ is a solution.
For $b = 1$ it is clearly the only solution (since $a \neq 1$).
For $b>1$, let $q : \mathbb{R}^{+} \to \mathbb{R}$ be the function defined by:
\begin{equation*}
 q\paren{y} =
 \begin{cases}
  \dfrac{y^b - a y + a - 1}{y - 1} =
  \dfrac{1 - y^b}{1 - y} - a
  &\text{for $y \neq 1$} \\
  b - a &\text{for $y = 1$}
 \end{cases}
\end{equation*}
which is well defined and continuous everywhere,
and satisfies the relation $y^b - a y + a - 1= \paren{y - 1} q\paren{y}$
for every $y \in \mathbb{R}^{+}$.

Claim: $q$ is strictly increasing. \newline
To prove this, first observe that $q$ is differentiable, with derivative:
\begin{equation*}
 q^{\prime}\paren{y} =
 \begin{cases}
  \dfrac{\paren{b-1} y^b - b y^{b-1} + 1}{\paren{y - 1}^2}
  &\text{for $y \neq 1$} \\
  \dfrac{1}{2} b \paren{b - 1} &\text{for $y = 1$}
 \end{cases}
\end{equation*}
It sufficies to show that $q^{\prime}$ is positive.
$q^{\prime}\paren{1} = \mfrac{1}{2} b \paren{b-1} > 0$,
so it remains to prove that $q^{\prime}\paren{y} > 0$ for $y \neq 1$.
For this, it is enough to prove that the numerator
$h\paren{y} = \paren{b-1} y^b - b y^{b-1} + 1$ is positive for $y \neq 1$
(because the denominator is $\paren{y - 1}^2$,
which is always positive for $y \neq 1$). \newline
The derivative $h^{\prime}\paren{y} = b\paren{b-1}\paren{y-1}y^{b-2}$
is negative for $0 < y < 1$,
positive for $y > 1$,
and is zero at $1$,
so $1$ is a point of global minimum for $h$;
thus $h\paren{y} > h\paren{1} = 0$ for $y \neq 1$, as required.

In particular, $q$ is strictly increasing in $\cinterval{1 - \mfrac{1}{a}}{1}$,
thus it admits at most one root in that range.
\begin{itemize}
  \item For $1 < b < a$,
        $q$ does not have any roots in the interval,
        because $q\paren{1} = b - a < 0$,
        so $q < 0$ on the whole interval (since it is increasing).
  \item For $b = a$,
        the root is $1$, since $q\paren{1} = b - a = 0$.
  \item For $b > a$,
        there is a root $\beta$ in the range,
        and it is neither $1 - \mfrac{1}{a}$ nor $1$.
        This is because $q\paren{1 - \mfrac{1}{a}} < 0$ and $q\paren{1} = b - a > 0$:
        \begin{equation*}
          q\paren{1 - \frac{1}{a}} =
          \frac{1 - \paren{1 - \mfrac{1}{a}}^b}{1 - 1 + \mfrac{1}{a}} - a =
          \frac{a^b - \paren{a - 1}^b}{a^{b-1}} - a =
          - \frac{\paren{a - 1}^b}{a^{b-1}} < 0
        \end{equation*}
\end{itemize}

So, in conclusion, in the interval $\cinterval{1 - \mfrac{1}{a}}{1}$:
\begin{itemize}
  \item for $\mfrac{1}{2} \leq r \leq \mfrac{1}{2} a$,
        the equation $y^{2 r} - a y + a - 1 = 0$
        admits only the solution $y = 1$;
  \item for $r > \mfrac{1}{2} a$,
        the equation $y^{2 r} - a y + a - 1 = 0$
        admits the two solutions $y = \beta$ and $y = 1$.
\end{itemize}

The solution $y = 1$ corresponds to $\cos\paren{x} = 0$,
that is $x = \mfrac{\pi}{2} + k \pi$ for $k \in \mathbb{Z}$.
Of those values of $x$,
the ones for which $k$ is even are indeed solutions to the original equation,
while those for which $k$ is odd are not, due to $\sin\paren{x}$ being negative in that case.
So this leaves the solutions
$x = \mfrac{\pi}{2} + 2 m \pi$ for $m \in \mathbb{Z}$.

The solution $y = \beta$ corresponds to $\cos\paren{x} = \pm \sqrt{a \paren{1 - \beta}}$.
Of the corresponding values of $x$, those for which $\sin\paren{x}$ is positive are solutions to the original equation,
while the others are not.
